\documentclass[12pt,a4paper]{article}

% ------------------------------------------------
% Pacotes
% ------------------------------------------------
\usepackage[utf8]{inputenc}
\usepackage[T1]{fontenc}
\usepackage[brazil]{babel}
\usepackage{graphicx}
\usepackage[table]{xcolor}
\usepackage{array}
\usepackage{tabularx}
\usepackage{geometry}
\usepackage{fancyhdr}
\usepackage{mathptmx} % Times New Roman
\usepackage{courier} % Adiciona suporte a negrito na fonte monospace
\usepackage{enumitem}\usepackage{amsfonts}
\usepackage{amsmath} % Recomendado para o comando \text{}
\usepackage{amssymb}
\usepackage{tikz}
\usetikzlibrary{positioning, shapes.geometric, arrows.meta}
\usepackage{listings}
\usepackage{multicol}

\lstset{
    language=Java,
    keywordstyle=\bfseries,
    showstringspaces=false,
    breaklines=true,
    frame=none,
    numbers=none,
    tabsize=2,
    commentstyle=\itshape\color{gray},
    morekeywords={this, class, private, public, String, System, out, println, Solo, Caminhao, Patio, Termometro},
    extendedchars=true,
    inputencoding=utf8
}

% TikZ Styles for ER Diagrams
\tikzset{
    block/.style={rectangle, draw, thick, minimum width=2cm, minimum height=1cm, fill=white},
    diamond_node/.style={diamond, draw, thick, minimum width=2cm, minimum height=1cm, fill=white, aspect=2},
    attribute/.style={ellipse, draw, thin, minimum width=1.5cm, minimum height=0.6cm, fill=white, font=\scriptsize},
    line/.style={draw, thick}, 
}
% ------------------------------------------------ 
% Configuração de página
% ------------------------------------------------
\geometry{
    a4paper,
    left=2cm,
    right=2cm,
    top=2cm,
    bottom=2cm,
    headheight=15pt % Adicione isso
}% Caixa para resposta discursiva
\newcommand{\resposta}[1]{
\noindent
\framebox[\textwidth]{
\begin{minipage}{0.98\textwidth}
\vspace*{#1}\null
\end{minipage}
}
}

\pagestyle{fancy}
\fancyhf{}
\renewcommand{\headrulewidth}{0pt}

\setlength{\parindent}{0pt}

% ------------------------------------------------
% Comandos auxiliares
% ------------------------------------------------
\newcolumntype{L}[1]{>{\raggedright\arraybackslash}p{#1}}
\newcolumntype{C}[1]{>{\centering\arraybackslash}p{#1}}

% ------------------------------------------------
% Documento
% ------------------------------------------------
\begin{document}

% ------------------------------------------------
% Cabeçalho
% ------------------------------------------------
\begin{tabularx}{\textwidth}{@{} L{4cm} X @{}}
    \raisebox{-0.5cm}{\includegraphics[width=4cm]{Universidade_de_Rio_Verde.jpg}} &
    \hfill
    \begin{tabular}{|>{\centering\arraybackslash}m{9.5cm}|}
        \hline
        \cellcolor{black}\color{white}\rule{0pt}{1.2cm}\textbf{LISTA DE REVISÃO -- N1} 
        \tabularnewline
        \hline
    \end{tabular} 
\end{tabularx}
\vspace{0.3cm}

% ------------------------------------------------  
% Tabela principal
% ------------------------------------------------
\begin{tabularx}{\textwidth}{|X|L{2cm}|}
\hline
\textbf{Disciplina} & \textbf{Campus} \\
\textit{Engenharia da Qualidade e Confiabilidade} & \textit{Rio Verde} \\
\hline
\textbf{Professor} & \textbf{Código} \\
\textit{Sergio Souza Novak} & \textit{ESW451} \\
\hline
\end{tabularx}
 
% ------------------------------------------------
% Barra de Atenção
% ------------------------------------------------
\vspace{0.2cm}

\begin{tabularx}{\textwidth}{|X|}
\hline
\cellcolor{black}\color{white}
\centering \textbf{ATENÇÃO:} Esta lista de revisão contém questões similares às da avaliação N1. Utilize-a como guia de estudos. \tabularnewline
\hline
\end{tabularx}

\vspace{0.3cm}

\textbf{ORIENTAÇÕES}  

\begin{itemize}
    \item Resolva todas as questões com atenção;
    \item Consulte o material de aula para verificar suas respostas;
    \item A interpretação de texto faz parte do exercício.
\end{itemize}

\vspace{0.5cm}

\begin{center}
    \rule{\textwidth}{0.4pt}
    \textbf{QUESTÕES}
\end{center}

\vspace{0.5cm}

% ================================================
% Questão 1 - AVAIL, POFOD, MTTF (cenário diferente)
% ================================================
\textbf{Questão 1} -- Um servidor web de uma empresa de comércio eletrônico é monitorado durante 6 meses. Nesse período, foram registradas 5 falhas, com os seguintes tempos de operação contínua entre elas: 400h, 600h, 200h, 350h e 450h. Após cada falha, o servidor levou, respectivamente, 8h, 12h, 4h, 6h e 10h para ser restaurado. Durante o mesmo período, o sistema recebeu 500 requisições críticas de pagamento, das quais 15 resultaram em falha.

Com base nesses dados e nos conceitos de métricas de confiabilidade, assinale a alternativa \textbf{correta}:

\vspace{0.3cm}

\begin{enumerate}[label=\textbf{(\Alph*)}]
  
  \item O MTTF (Mean Time To Failure) é calculado somando-se os tempos de reparo e dividindo pelo número de falhas, resultando em 8 horas.

  \item O POFOD (Probability of Failure on Demand) é calculado dividindo-se o número de falhas (15) pelo número total de requisições críticas (500), resultando em 0,03 (ou 3\%).

  \item O AVAIL (Disponibilidade) é calculado dividindo-se o tempo total de reparo pelo tempo total de operação, resultando em aproximadamente 0,02.

  \item O MTTF é obtido dividindo-se o número de falhas pelo tempo total de operação, resultando em 0,0025 falhas por hora.

  \item O AVAIL e o MTTF representam a mesma métrica expressa em unidades diferentes: o AVAIL em percentual e o MTTF em horas.
\end{enumerate}

\vspace{0.5cm}

% ================================================
% Questão 2 - Defeito, Erro, Estado de Erro, Estado de Falha (cenário diferente)
% ================================================
\textbf{Questão 2} -- Uma equipe de desenvolvimento está criando um sistema de navegação GPS para veículos autônomos. Durante a implementação do módulo de cálculo de distância entre dois pontos geográficos, um programador utiliza a fórmula de distância euclidiana plana em vez da fórmula de Haversine (que considera a curvatura da Terra). Em trajetos curtos dentro de uma cidade, a diferença entre os valores é desprezível e o veículo segue normalmente. Porém, ao calcular uma rota interestadual de 800 km, o sistema estima a distância como 785 km, fazendo com que o veículo tente seguir por uma rota alternativa mais curta que passa por uma estrada interditada, resultando na imobilização do veículo.

\vspace{0.2cm}

Com base no cenário descrito e nos conceitos de \textbf{defeito}, \textbf{erro}, \textbf{estado de erro} e \textbf{estado de falha}, responda:

\begin{enumerate}[label=\textbf{\alph*)}]
    \item Identifique e defina cada um dos quatro conceitos (\textit{defeito}, \textit{erro}, \textit{estado de erro} e \textit{estado de falha}) aplicando-os ao cenário apresentado.
    \item Explique por que o defeito não se manifestou imediatamente nos trajetos curtos, mas levou a uma falha nos trajetos longos.
\end{enumerate}

\vspace{0.2cm}
\resposta{7cm}

\vspace{0.5cm}

% ================================================
% Questão 3 - Programação N-Versões (cenário diferente)
% ================================================
\textbf{Questão 3} -- Um hospital de referência decide implementar um sistema de dosagem automática de medicamentos para pacientes em UTI. Devido à criticidade do sistema, a equipe de engenharia adota a técnica de \textbf{Programação N-Versões}: quatro equipes independentes desenvolvem, a partir da mesma especificação, implementações distintas do algoritmo de cálculo de dosagem. Em tempo de execução, as quatro versões são executadas simultaneamente e um módulo comparador seleciona o resultado de consenso majoritário.

\vspace{0.2cm}

Sobre a técnica de Programação N-Versões e sua aplicação no cenário descrito, assinale a alternativa \textbf{correta}:

\vspace{0.3cm}

\begin{enumerate}[label=\textbf{(\Alph*)}]
    \item A técnica elimina completamente a possibilidade de falhas no sistema, pois, com quatro versões independentes, é estatisticamente impossível que mais de duas produzam o mesmo erro.

    \item O módulo comparador (árbitro) é um componente trivial e não precisa ser testado com rigor, pois sua função é apenas selecionar a saída mais frequente entre as versões.

    \item Uma limitação importante da técnica é que ela não protege contra erros presentes na própria especificação, pois todas as versões são desenvolvidas a partir do mesmo documento de requisitos; se a especificação estiver incorreta, todas as versões poderão produzir o mesmo resultado errado.

    \item A Programação N-Versões é economicamente vantajosa porque o custo total de desenvolvimento é equivalente ao de uma única versão, já que as equipes compartilham bibliotecas e frameworks.

    \item A técnica funciona de forma idêntica à redundância de hardware (\\textit{hot standby}), na qual cópias idênticas do mesmo software são executadas em processadores diferentes.
\end{enumerate}

\vspace{0.5cm}

% ================================================
% Questão 4 - Apache Hadoop (cenário diferente)
% ================================================
\textbf{Questão 4} -- Uma rede de varejo com mais de 2.000 lojas em todo o Brasil deseja analisar o comportamento de compra de seus clientes a partir dos dados de transações de cartões de fidelidade. São gerados aproximadamente 50 milhões de registros diários, totalizando vários terabytes por mês. A equipe de dados propõe o uso do \textbf{Apache Hadoop} para viabilizar essa análise.

\vspace{0.2cm}

Considerando os conceitos fundamentais do Apache Hadoop, assinale a alternativa \textbf{correta}:

\vspace{0.3cm}

\begin{enumerate}[label=\textbf{(\Alph*)}]
    \item O HDFS (Hadoop Distributed File System) armazena cada arquivo em um único nó do cluster para maximizar a velocidade de leitura sequencial, sem replicação, pois a redundância é garantida pelo hardware RAID dos servidores.

    \item O modelo MapReduce divide o processamento em duas fases principais: a fase \textit{Map}, que aplica uma função de transformação aos dados de entrada em paralelo nos nós do cluster, e a fase \textit{Reduce}, que agrega os resultados intermediários para produzir a saída final. Esse modelo permite o processamento eficiente de grandes volumes de dados de forma distribuída e tolerante a falhas.

    \item O Hadoop exige que todos os dados estejam estruturados em formato tabular (como em bancos de dados relacionais) antes de serem carregados no HDFS, pois o MapReduce só opera sobre dados com esquema fixo definido previamente.

    \item O Hadoop substitui completamente bancos de dados tradicionais, pois oferece suporte nativo a consultas SQL em tempo real com latência inferior a milissegundos, sendo ideal para sistemas transacionais de alta frequência.

    \item O processamento no Hadoop é centralizado em um único nó mestre que executa toda a lógica de negócio, enquanto os nós escravos servem apenas como unidades de armazenamento passivo de dados.
\end{enumerate}

\vspace{0.5cm}

% ================================================
% Questão 5 - Escalonamento Vertical e Horizontal (V/F)
% ================================================
\textbf{Questão 5} -- Sobre as estratégias de \textbf{escalonamento vertical} (\textit{scale up}) e \textbf{escalonamento horizontal} (\textit{scale out}) de sistemas, julgue as afirmativas abaixo como \textbf{(V) Verdadeiro} ou \textbf{(F) Falso}:

\vspace{0.3cm}

\begin{enumerate}[label=\textbf{(\hspace{0.6cm})} ,leftmargin=2cm]

    \item O escalonamento vertical é geralmente mais simples de implementar inicialmente, pois não exige mudanças na arquitetura da aplicação — basta atualizar o hardware do servidor existente.

    \item Uma desvantagem do escalonamento horizontal é a necessidade de mecanismos de sincronização e consistência de dados entre os nós, o que aumenta a complexidade do sistema.

    \item O escalonamento vertical oferece melhor tolerância a falhas do que o horizontal, pois concentra todos os recursos em uma única máquina robusta, eliminando pontos de falha distribuídos.

    \item Sistemas como o Apache Hadoop e o Apache Kafka são exemplos de tecnologias projetadas para operar com escalonamento horizontal, distribuindo dados e processamento entre múltiplos nós.

    \item O escalonamento horizontal é especialmente adequado para aplicações \textit{stateless} (sem estado), pois qualquer nó pode atender qualquer requisição sem depender de dados de sessão armazenados localmente.

    \item O escalonamento vertical não possui nenhuma limitação prática, pois é sempre possível adicionar mais recursos (CPU, RAM, disco) a um único servidor de forma ilimitada.

\end{enumerate}

\vspace{0.5cm}

% ================================================
% Questão 6 - Arquitetura do Sistema de Proteção (Sommerville) - Discursiva
% ================================================
\textbf{Questão 6} -- Considere um sistema de controle de uma usina nuclear que possui dois subsistemas: um \textbf{sistema de controle} responsável por regular a temperatura do reator e um \textbf{sistema de proteção} responsável por desligar automaticamente o reator caso a temperatura ultrapasse um limite crítico. Ambos possuem seus próprios sensores e lógica de decisão independentes.

\vspace{0.2cm}

Com base nos conceitos de Sommerville sobre \textbf{sistemas críticos de segurança} e na arquitetura de sistemas de proteção, responda:

\begin{enumerate}[label=\textbf{\alph*)}]
    \item Por que é fundamental que o sistema de proteção possua \textbf{sensores próprios}, independentes dos sensores utilizados pelo sistema de controle?
    \item Explique o que poderia acontecer se o sistema de proteção compartilhasse o mesmo software e hardware do sistema de controle. Qual princípio de confiabilidade seria violado?
\end{enumerate}

\vspace{0.2cm}
\resposta{6cm}

\vspace{0.5cm}

% ================================================
% Questão 7 - ROCOF (Rate of Occurrence of Failures)
% ================================================
\textbf{Questão 7} -- Um sistema de gerenciamento de estoque em uma fábrica é monitorado ao longo de quatro trimestres consecutivos. No primeiro trimestre (2.160h de operação), foram registradas 6 falhas; no segundo trimestre (2.160h), 4 falhas; no terceiro trimestre (2.160h), 8 falhas; e no quarto trimestre (2.160h), 2 falhas.

\vspace{0.2cm}

Com base nesses dados, responda:

\begin{enumerate}[label=\textbf{\alph*)}]
    \item Calcule o \textbf{ROCOF} (Rate of Occurrence of Failures) para cada trimestre e para o período total.
    \item O sistema está se tornando mais confiável ou menos confiável ao longo do tempo? Justifique sua resposta com base nos valores de ROCOF calculados.
    \item Se a empresa deseja que o ROCOF total do próximo ano seja inferior a 0,002 falhas/hora, quantas falhas, \textbf{no máximo}, o sistema poderia apresentar ao longo de 8.640h de operação?
\end{enumerate}

\vspace{0.2cm}
\resposta{6cm}

\vspace{0.5cm}

% ================================================
% Questão 8 - Tolerância a Falhas: Recovery Blocks vs N-Versões
% ================================================
\textbf{Questão 8} -- Uma empresa de energia eólica está desenvolvendo o software de controle de suas turbinas. O sistema deve ajustar automaticamente o ângulo das pás com base na velocidade do vento. Dois engenheiros propõem abordagens diferentes para garantir a tolerância a falhas do módulo de ajuste:

\begin{itemize}
    \item \textbf{Engenheiro A} propõe utilizar \textbf{Recovery Blocks} (Blocos de Recuperação): uma versão primária do algoritmo é executada e seu resultado é verificado por um teste de aceitação; caso o teste falhe, uma versão alternativa (backup) é executada.
    \item \textbf{Engenheiro B} propõe utilizar \textbf{Programação N-Versões}: três versões independentes do algoritmo são executadas simultaneamente e o resultado é decidido por votação majoritária.
\end{itemize}

\vspace{0.2cm}

Sobre as duas abordagens, assinale a alternativa \textbf{correta}:

\vspace{0.3cm}

\begin{enumerate}[label=\textbf{(\Alph*)}]
    \item Na abordagem de Recovery Blocks, todas as versões são executadas simultaneamente, assim como na Programação N-Versões, diferindo apenas no critério de decisão: teste de aceitação contra votação.

    \item A Programação N-Versões tem como vantagem a detecção mais rápida de falhas, pois o resultado é obtido pela comparação simultânea de múltiplas saídas, sem a necessidade de reexecução. Já os Recovery Blocks podem apresentar maior latência em caso de falha da versão primária, pois precisam executar sequencialmente as versões alternativas.

    \item Os Recovery Blocks são superiores à Programação N-Versões em todos os critérios, pois exigem menos recursos computacionais, são mais simples de implementar e oferecem a mesma garantia de tolerância a falhas.

    \item Ambas as técnicas garantem que o sistema nunca falhará, desde que ao menos duas versões tenham sido desenvolvidas por equipes independentes.

    \item Os Recovery Blocks não necessitam de um teste de aceitação, pois a versão de backup é executada automaticamente após qualquer saída da versão primária, independentemente de seu valor.
\end{enumerate}

\vspace{0.5cm}

% ================================================
% Questão 9 - Conceitos Gerais de Confiabilidade
% ================================================
\textbf{Questão 9} -- Relacione os conceitos da coluna da esquerda com suas respectivas definições na coluna da direita:

\vspace{0.3cm}

\begin{tabularx}{\textwidth}{|c|L{3.2cm}|c|X|}
\hline
\textbf{N°} & \textbf{Conceito} & \textbf{(\ \ )} & \textbf{Definição} \\
\hline
1 & Confiabilidade (\textit{Reliability}) & & Probabilidade de o sistema causar dano a pessoas ou ao ambiente quando uma falha ocorre. \\
\hline
2 & Disponibilidade (\textit{Availability}) & & Probabilidade de o sistema estar operacional e acessível quando solicitado. \\
\hline
3 & Segurança (\textit{Safety}) & & Capacidade do sistema de funcionar sem falhas em condições normais durante um período de tempo especificado. \\
\hline
4 & Proteção (\textit{Security}) & & Capacidade do sistema de se proteger contra ataques externos, acidentais ou deliberados. \\
\hline
\end{tabularx}

\vspace{0.2cm}

A sequência correta, de cima para baixo, é:

\vspace{0.2cm}

\begin{enumerate}[label=\textbf{(\Alph*)}]
    \item 3, 2, 1, 4
    \item 4, 2, 3, 1
    \item 1, 3, 2, 4
    \item 2, 1, 4, 3
    \item 3, 1, 2, 4
\end{enumerate}

\vspace{0.5cm}

% ================================================
% Questão 10 - Questão Integradora (Discursiva)
% ================================================
\textbf{Questão 10} -- Uma fintech (empresa de tecnologia financeira) está projetando um sistema de processamento de transações PIX que deve operar 24 horas por dia, 7 dias por semana. O sistema precisa atender a milhões de transações diárias com tempo de resposta inferior a 2 segundos. A equipe de arquitetura está debatendo as seguintes decisões:

\begin{itemize}
    \item Utilizar escalonamento horizontal com múltiplos nós de processamento;
    \item Implementar redundância de software com técnicas de tolerância a falhas;
    \item Definir métricas de confiabilidade para monitoramento contínuo.
\end{itemize}

\vspace{0.2cm}

Com base nos conceitos estudados, responda:

\begin{enumerate}[label=\textbf{\alph*)}]
    \item Justifique por que o \textbf{escalonamento horizontal} é mais adequado do que o vertical para este cenário. Cite pelo menos dois argumentos.
    \item Indique quais \textbf{métricas de confiabilidade} (dentre AVAIL, POFOD, ROCOF e MTTF) seriam mais relevantes para monitorar este sistema e justifique sua escolha.
    \item Proponha uma estratégia de \textbf{tolerância a falhas} (N-Versões ou Recovery Blocks) para o módulo de validação de transações e justifique sua escolha.
\end{enumerate}

\vspace{0.2cm}
\resposta{8cm}

\end{document}
